\section{Odd scalar channel and axial stress tensor}\label{sec:axial-channel}

\subsection{Diagonal angular momentum and parity}

The spherical metric does not imply that every scalar perturbation is polar.
Spatial rotations act simultaneously on the orbital harmonics and on the
internal multiplet. It is therefore convenient to make a unitary change of
internal basis and couple an orbital representation $L$ to the background
representation $\ell$. With $s=-\ell,\ldots,\ell$, define the normalized
background angular vector and coupled perturbation harmonic
\begin{align}
B_s^{(\ell)}&=\frac{(-1)^\ell}{\sqrt{2\ell+1}}Y_{\ell s}^{*},
\label{eq:background-angular-vector}\\
P_s^{JML}&=\sum_{m_L=-L}^{L}
\langle Lm_L,\ell s\vert JM\rangle Y_{Lm_L}.
\label{eq:coupled-harmonic}
\end{align}
The radial field multiplying Eq.~\eqref{eq:background-angular-vector} is
$F=\sqrt{2\ell+1}\,\psi_\ell$, so this normalized basis leaves the physical
sum over fields unchanged. The allowed values satisfy
$|J-\ell|\leq L\leq J+\ell$. For $\ell=1$, the $L=J-1,J+1$ channels are even
and the $L=J$ channel is odd \cite{ChavezNambo2024}.

The scalar density bilinear relevant at leading nonrelativistic order is
\begin{equation}
\mathcal D^{JML}=\sum_s B_s^{*}P_s^{JML}.
\label{eq:density-bilinear}
\end{equation}
Its reduced coefficient is proportional to
$\langle L0,\ell0\vert J0\rangle$. For $\ell=1$ and $J=L=2$, the sum
$J+L+\ell=5$ is odd; hence
\begin{equation}
\mathcal D^{2M2}=0
\qquad (M=-2,-1,0,1,2).
\label{eq:density-null}
\end{equation}
This is the group-theoretic origin of the density-silent odd sector in the
Schr\"odinger--Poisson limit. It does not imply that every component of the
relativistic stress tensor vanishes.

\subsection{Axial vector and tensor projections}

Let $D_A$ and $\varepsilon_{AB}$ be the covariant derivative and volume form
of the unit two-sphere. We use the normalized axial vector harmonic and its
symmetrized derivative,
\begin{equation}
X_A^{JM}=-\frac{\varepsilon_A{}^{B}D_BY_{JM}}{\sqrt{J(J+1)}},
\qquad
X_{AB}^{JM}=D_{(A}X_{B)}^{JM}.
\label{eq:axial-harmonics}
\end{equation}
Our normalization gives
\begin{equation}
\int X_A^{*}X^A\,\dd\Omega=1,
\qquad
\int X_{AB}^{*}X^{AB}\,\dd\Omega
=\frac{(J-1)(J+2)}{2}.
\label{eq:axial-norms}
\end{equation}
The angular structures entering the mixed and trace-free angular components
of the scalar stress tensor are
\begin{equation}
S_A^{JML}=\sum_s B_s^{*}D_AP_s^{JML},
\qquad
Q_{AB}^{JML}=\sum_s D_{(A}B_s^{*}D_{B)}P_s^{JML}.
\label{eq:axial-bilinears}
\end{equation}
Direct evaluation for $\ell=1$, $J=L=2$ yields the pointwise identities
\begin{equation}
S_A^{2M2}=-\frac{\ii}{2\sqrt{\pi}}X_A^{2M},
\qquad
Q_{AB}^{2M2}=\frac{\ii}{2\sqrt{\pi}}X_{AB}^{2M}.
\label{eq:axial-pointwise}
\end{equation}
Consequently,
\begin{equation}
\int X^{A*}S_A\,\dd\Omega=-\frac{\ii}{2\sqrt{\pi}},
\qquad
\int X^{AB*}Q_{AB}\,\dd\Omega=\frac{\ii}{\sqrt{\pi}}.
\label{eq:axial-projections}
\end{equation}
The factor of two between the second pointwise and integrated coefficients is
the $J=2$ tensor-harmonic norm in Eq.~\eqref{eq:axial-norms}. Overall phases
change if the tensor-harmonic basis is rephased, but vanishing, magnitudes,
and $M$ degeneracy do not.

\begin{table}[t]
\centering
\caption{Independent quadrature check of the $\ell=1$, $J=L=2$ angular
identities. The exact values are
$|\langle X,S\rangle|=1/(2\sqrt\pi)=0.2820947917738781$ and
$|\langle X_{(2)},Q\rangle|=1/\sqrt\pi=0.5641895835477563$.}
\label{tab:axial-projection}
\small
\begin{tabular}{@{}rccc@{}}
\toprule
$M$ & $\|\mathcal D\|_2$ & $\operatorname{Im}\langle X,S\rangle$
& $\operatorname{Im}\langle X_{(2)},Q\rangle$ \\
\midrule
$-2$ & $2.83\times10^{-17}$ & $-0.2820947917738770$ & $0.5641895835477598$ \\
$-1$ & $2.68\times10^{-17}$ & $-0.2820947917738799$ & $0.5641895835477542$ \\
$0$  & $1.01\times10^{-17}$ & $-0.2820947917738773$ & $0.5641895835477541$ \\
$1$  & $2.69\times10^{-17}$ & $-0.2820947917738799$ & $0.5641895835477542$ \\
$2$  & $2.83\times10^{-17}$ & $-0.2820947917738770$ & $0.5641895835477598$ \\
\bottomrule
\end{tabular}
\end{table}

The $M=0$ result was also obtained symbolically from explicit low-order
spherical harmonics, without numerical quadrature. The numerical calculation
uses Gauss--Legendre integration in $\cos\theta$ and uniform Fourier nodes in
$\phi$. Increasing the grid from $24\times32$ to $48\times64$ changes each
projected coefficient by less than $2\times10^{-12}$; at the reported
resolution the spread over $M$ is $2.9\times10^{-15}$ for the vector and
$5.8\times10^{-15}$ for the tensor projection.

\subsection{Linearized relativistic source}

To expose the physical content, consider one frequency sideband in the same
time convention as the background,
\begin{equation}
\Phi_s^{(0)}=e^{\ii\omega t}F(r)B_s,
\qquad
\delta\Phi_s^{(+)}=e^{\ii\omega t-\ii\sigma t}u(r)P_s^{2M2}.
\label{eq:axial-sideband}
\end{equation}
The companion sideband required by the full complex perturbation contributes
the conjugate angular structure with an independent radial amplitude. For the
sideband in Eq.~\eqref{eq:axial-sideband}, Eq.~\eqref{eq:density-null} implies
$\sum_sP_sD_AB_s^*=-S_A$. Substitution into the stress tensor
\eqref{eq:stress} gives
\begin{align}
\delta T_{tA}^{(+,\mathrm{scalar})}
&=-\frac{\ii}{2}F(2\omega-\sigma)uS_A,
\label{eq:stress-ta-pre}\\
\delta T_{rA}^{(+,\mathrm{scalar})}
&=\frac{1}{2}(F'u-Fu')S_A,
\label{eq:stress-ra-pre}\\
\delta T_{AB}^{(+,\mathrm{scalar,odd})}
&=FuQ_{AB}.
\label{eq:stress-ab-pre}
\end{align}
Using Eq.~\eqref{eq:axial-pointwise}, the projected sources are therefore
\begin{align}
\delta T_{tA}^{(+,\mathrm{scalar})}
&=-\frac{F(2\omega-\sigma)u}{4\sqrt\pi}X_A^{2M},
\label{eq:stress-ta}\\
\delta T_{rA}^{(+,\mathrm{scalar})}
&=-\frac{\ii(F'u-Fu')}{4\sqrt\pi}X_A^{2M},
\label{eq:stress-ra}\\
\delta T_{AB}^{(+,\mathrm{scalar,odd})}
&=\frac{\ii Fu}{2\sqrt\pi}X_{AB}^{2M}.
\label{eq:stress-ab}
\end{align}
These equations demonstrate that the cancellation of the density channel does
not remove the relativistic angular current or odd anisotropic stress.

For completeness, write the odd metric perturbation as
\begin{equation}
h_{tA}=h_0X_A^{2M},\qquad h_{rA}=h_1X_A^{2M},\qquad
h_{AB}=h_2X_{AB}^{2M}.
\label{eq:axial-metric}
\end{equation}
If
\begin{equation}
\mathcal K_0=\sum_s\left(
\nabla^c\Phi_s^{(0)*}\nabla_c\Phi_s^{(0)}
+\mu^2|\Phi_s^{(0)}|^2\right),
\end{equation}
then variation of the explicit metric in Eq.~\eqref{eq:stress} adds
$-\mathcal K_0h_{ab}/2$. Terms proportional to the background
$g_{AB}$ are pure angular trace and have zero axial-tensor projection. Thus
the complete axial projections for this sideband are obtained from
Eqs.~\eqref{eq:stress-ta}--\eqref{eq:stress-ab} by adding
$-\mathcal K_0h_0X_A/2$, $-\mathcal K_0h_1X_A/2$, and
$-\mathcal K_0h_2X_{AB}/2$, respectively. Regge--Wheeler gauge sets $h_2=0$
for $J\geq2$. The angular system therefore closes on the usual axial metric
harmonics, but with a genuine scalar source.

\subsection{Null control and weak-field interpretation}

For an ordinary $\ell=0$ background, $B=Y_{00}$ is constant. The vector
bilinear is then proportional to $D_AY_{JM}$ and is purely polar, while
$Q_{AB}=0$. Both axial projections vanish exactly. This control shows that
Eq.~\eqref{eq:axial-pointwise} is not an artifact of projecting an arbitrary
scalar gradient: it requires the internal orbital structure of the
$\ell=1$ background.

In the nonrelativistic limit the leading mass-density perturbation is
proportional to Eq.~\eqref{eq:density-bilinear}, so the odd sector does not
source the Schr\"odinger--Poisson potential, in agreement with
Ref.~\cite{ChavezNambo2024}. The relativistic current and anisotropic-stress
terms are nevertheless nonzero. Their metric response becomes
post-Newtonianly suppressed as compactness decreases. The radial axial
equations, massive-scalar sheets, and outgoing gravitational boundary
conditions developed below quantify that response and resolve both the QNM
frequency and its damping law.

\subsection{Verified metric--scalar radial coupling}

The first dynamical step beyond the angular proof is to vary the
Klein--Gordon operator with respect to the odd metric.  In Regge--Wheeler
gauge, let the positive sideband have radial amplitude $u_+(r)$ and total
time dependence $e^{\ii(\omega-\sigma)t}$.  Direct coordinate variation of
$\Box\Phi$, without using the Einstein equations or a published component
formula, gives
\begin{align}
0={}&\frac{1}{\alpha\gamma r^2}\frac{\dd}{\dd r}
\left(\frac{\alpha r^2}{\gamma}u_+'\right)
+\left[\frac{(\omega-\sigma)^2}{\alpha^2}-\mu^2-
\frac{L(L+1)}{r^2}\right]u_+ \nonumber\\
&+\frac{\kappa_V}{r^2}\left[
\frac{\ii F(2\omega-\sigma)}{\alpha^2}h_0
-\frac{1}{\gamma^2}\left\{Fh_1'
+\left(2F'+F\left(\frac{\alpha'}{\alpha}-
\frac{\gamma'}{\gamma}\right)\right)h_1\right\}\right],
\label{eq:kg-positive-sideband}
\end{align}
where $L=2$ and $\kappa_V=-\ii/(2\sqrt\pi)$ in the convention of
Eq.~\eqref{eq:axial-pointwise}.  The companion amplitude $u_-$ obeys the
same radial differential operator with $(\omega-\sigma)^2$ replaced by
$(\omega+\sigma)^2$; its time-current coupling changes to
$-\ii F(2\omega+\sigma)h_0/\alpha^2$, while the radial $h_1$ coupling is
unchanged in this phase convention.

Equation~\eqref{eq:kg-positive-sideband} has been checked by two algebraically
independent calculations.  The first varies
$g^{ab}(\partial_a\partial_b-\Gamma^c{}_{ab})\Phi$ in coordinates on a
general spherical background.  The second expands the compact expression
shown above.  Their symbolic difference simplifies identically to zero.
Setting $h_0=h_1=0$ recovers the massive scalar sideband equation; setting
$F=0$ removes the scalar--metric coupling.

This calculation also corrects a useful bookkeeping misconception.  The two
scalar sidebands obey second-order radial equations.  Therefore a local
first-order formulation cannot close on only $(h_0,h_1,u_+,u_-)$.  Before
Einstein constraints are used to eliminate variables, the minimal state is
six-dimensional, for example
\begin{equation}
\bm{Y}=(h_0,h_1,u_+,u_+',u_-,u_-')^{T}.
\label{eq:minimal-axial-state}
\end{equation}
The next section combines Eq.~\eqref{eq:kg-positive-sideband} and its
companion with two axial Einstein equations, imposes regular central data and
outgoing/decaying channel conditions, and performs a declared complex search.
It also records which certification gates are closed and which remain open,
so the reported frequency is interpreted as a strong candidate pole rather
than a completed independent damping measurement.
